\section{Design implications}
\label{app:design}

The mathematics above answers a question a system designer actually has: given a
budget on what survives an episode, what must be kept so that someone can still
be paid, credited or blamed for their own work. Four rules follow, each traceable
to a specific result rather than to intuition.

\paragraph{R1. Preserve agent identity before preserving content.}
This is the sharpest empirical finding of Section~\ref{app:episodes} and it is
not what a compression-minded engineer expects. The record that keeps every
artifact in full and discards only \emph{who produced each one} retains more of
the trace than a single bit per agent, and has less than a third of the
accountability margin ($0.1425$ against $0.5000$ at $n=5$; at $n = 3$ it has
\emph{more} labels and still less than half the margin). Attribution is what is
being bought, and content is not a substitute for it. The magnitude varies: across
$n = 3,\dots,7$ the cost of removing identity, measured against the storage it
saves, is reliably greater than one but usually between $1.4$ and $2.3$ rather
than the near-sixfold figure that case gives. The direction is the rule; the
size is not.

\paragraph{R2. Keep one witness per agent, for every agent.}
One retained bit per agent is flat with margin $1$: the tension does not exist
there. Under the episode model it is not even a compromise, matching the full
trace on both reconstruction and margin while using $32$ labels instead of $243$
(Table~\ref{tab:episodes}). The important word is \emph{every}: a record that
keeps per-agent bits for a fixed subset and drops the rest has margin exactly
zero, because the margin is a minimum over agents and the dropped ones are
invisible (Table~\ref{tab:mechanisms}). Partial per-agent retention on a schedule
fixed in advance buys nothing at all.

\paragraph{R3. Do not settle from a global count.}
An aggregate such as ``how many subtasks completed'' has accountability margin
exactly $1/n$ --- it decays with the size of the team, which is precisely the
regime where deferred settlement matters most. Note this is weaker than the
single-deviation intuition suggests: one agent always moves the count by one, yet
mixtures of teammate behaviour cancel the intermediate shifts
(Section~\ref{sec:experiment}). Counts are cheap and feel informative; they are
close to worthless for attribution at scale.

\paragraph{R4. More retained content does not mean more accountability.}
The two axes are close to independent, and the paper's central object is a pair
of records that makes this concrete: at $n = 5$ two ten-label records incur
\emph{identical} distortion $\Dnull(10)$, and one supports margin $1/2$ while the
other supports none at any price. A designer comparing candidate records on
reconstruction quality alone---which is the natural thing to measure---cannot
tell them apart. Accountability has to be measured separately, and
Theorem~\ref{thm:elim} makes that measurement cheap: one linear program per
agent, no search over payment rules.

\paragraph{A caveat that cuts the other way.}
Two of these rules are robustness claims, and Section~\ref{sec:experiment} found
one boundary. Selection that \emph{depends on what agents did} is much better
than selection fixed in advance: retaining bits for the top half of agents by
observed quality keeps $0.2624$ of the available $0.5000$, not the zero that R2's
fixed-subset case gives, because an agent who shirks changes its own probability
of being recorded. If a budget forces dropping agents, drop them on the basis of
the episode rather than on a schedule.
